\section{결론}

본 논문은 멀티-툴 도구 선택을 위한 훈련-없는 \emph{부호 반영 회전 스티어링}을 제시했다. 카탈로그에서 직접 유도한 per-head 온톨로지 부분공간 위에서 질의를 $Q' = Q + \beta Q P_{\mathrm{ont}}$로 회전시키되, 회전의 부호 $\beta$는 도메인마다 다르다는 사실을 네 Qwen 도메인(MetaTool Subtask4, $\tau^2$-retail/telecom/airline)에서 locked 수치로 확인했다 — Telecom에서 $\beta>0$이 $+24.78$pp, MetaTool과 Retail에서 $\beta<0$이 각각 $+2.28$pp, $+5.11$pp. 핵심 기여는 이 부호를 baseline forward pass 한 번에 예측하는 정리~\ref{thm:beta-star}이며, softmax gradient identity에서 $\operatorname{sign}(\beta^*) = \operatorname{sign}(\bar r_{\mathcal{G}} - \bar r)$로 직접 유도된다. Adaptive Sign Routing 알고리즘은 이 예측기를 추론 시점에 실행하여 부호 선택을 hyperparameter search가 아닌 \emph{이론적 결정}으로 만든다. 정리~\ref{thm:no-memory}은 가장 가까운 선행 방법 SEKA와 AdaSEKA가 정지 key-side 증폭 형식 때문에 방출 상태를 구조적으로 담지 못함을 보이며, 그 footprint는 \texttt{repeated\_first\_tool\_rate}의 상승으로 데이터에 남는다. 효율성 측면에서 본 방법은 파인튜닝 대비 학습 비용 0, 프롬프트 스캐폴딩 대비 토큰 오버헤드 0, 저장 $\approx 2.6$ MB로 추론 시 저랭크 hook이며, KV-cache 최적화와 직교 결합 가능하다. 부호가 도메인 의존적이라는 사실과 그 부호를 예측하는 정리가 함께 제시된 것은 정지 활성 스티어링 문헌 전체에 ``부호 선택''이라는 새 축을 추가한다.
